\SourceLeaf{279}%
Durchführung derselben mit zehnziffrigen Logarithmen, vollkommen zureiche. Jeder, der diese Tafel zur Berechnung von Messungen innerhalb dieser Zone benutzen will, wird, wenn eine geringere Schärfe ihm genügt (und diess ist allerdings der gewöhnlichste Fall) nach Gefallen einige der letzten Decimalen weglassen. In welcher Form man übrigens auch die \textit{Resultate} einer Messung darstellen mag, so sollte diess, consequenter Weise, immer in einer Schärfe geschehen, die der Schärfe der Messungen selbst entsprechend ist, so dass man aus den Zahlen der Resultate immer rückwärts die beobachteten Grössen eben so scharf wieder finden kann, wie sie gemessen waren. Wählt man also dazu ausschliesslich die Längen und Breiten, so würde trigonometrischen Messungen selbst von nur mässiger Schärfe, durchaus nicht ihr Recht widerfahren, wenn man die Resultate nur in solcher Schärfe ansetzen wollte, wie Längen und Breiten sich auf astronomischem Wege bestimmen lassen: man würde dadurch nur einen falschen Maassstab für die Güte der Arbeit erhalten, und sich oft gerade der durchgreifendsten Prüfungen dieser Güte entäussern.
\Article{11}
\indent Die Benutzung der hier betrachteten conformen Übertragung der Ellipsoidfläche auf die Kugelfläche zur Berechnung trigonometrischer Messungen kann auf mehr als Eine Art geschehen: in der gegenwärtigen Abhandlung wird nur von der unmittelbaren Benutzung die Rede sein; andere abgeleitete Arten, sie zu jenem Zwecke zu benutzen, sollen einer zweiten Abhandlung vorbehalten bleiben.

Die unmittelbare Benutzung ist im Wesentlichen schon in der oben angeführten Schrift kurz angedeutet. Ein auf der Oberfläche des Ellipsoids durch kürzeste oder sogenannte geodaetische Linien gebildetes System von Dreiecken wird auf der Oberfläche der Kugel durch ein Dreieckssystem dargestellt, worin die Winkel den entsprechenden auf dem Sphaeroid genau gleich sind, die Seiten hingegen, wenn sie nicht Meridianbögen sind, zwar nicht in aller Strenge Bögen Grösster Kreise werden, aber doch von solchen so wenig abweichen, dass sie in den meisten Fällen als damit ganz zusammenfallend betrachtet werden dürfen, oder dass wenigstens die Abweichung, da, wo die grösste Genauigkeit gefordert wird, mit aller nöthigen Schärfe leicht berechnet werden kann, immer vorausgesetzt, dass
